%----------%
% Theorems %
%----------%

\usepackage{thmtools}
\usepackage[framemethod=TikZ]{mdframed}
\mdfsetup{skipabove=1em,skipbelow=0em}

\theoremstyle{definition}

\declaretheoremstyle[
  headfont=\bfseries\sffamily\color{ForestGreen!70!black}, bodyfont=\normalfont,
  mdframed={
    linewidth=2pt,
    rightline=false, topline=false, bottomline=false,
    linecolor=ForestGreen, backgroundcolor=ForestGreen!5,
  }
]{thmgreenbox}

\declaretheoremstyle[
  headfont=\bfseries\sffamily\color{NavyBlue!70!black}, bodyfont=\normalfont,
  mdframed={
    linewidth=2pt,
    rightline=false, topline=false, bottomline=false,
    linecolor=NavyBlue, backgroundcolor=NavyBlue!5,
  }
]{thmbluebox}

\declaretheoremstyle[
  headfont=\bfseries\sffamily\color{NavyBlue!70!black}, bodyfont=\normalfont,
  mdframed={
    linewidth=2pt,
    rightline=false, topline=false, bottomline=false,
    linecolor=NavyBlue
  }
]{thmblueline}

\declaretheoremstyle[
  headfont=\bfseries\sffamily\color{RawSienna!70!black}, bodyfont=\normalfont,
  mdframed={
    linewidth=2pt,
    rightline=false, topline=false, bottomline=false,
    linecolor=RawSienna, backgroundcolor=RawSienna!5,
  }
]{thmredbox}

\declaretheoremstyle[
headfont=\bfseries\sffamily\color{RawSienna!70!black}, bodyfont=\normalfont,
mdframed={
  linewidth=2pt,
  rightline=false, leftline=false, topline=true, bottomline=true,
  linecolor=RawSienna,
}
]{thmredlinesupdown}

\declaretheoremstyle[
headfont=\bfseries\sffamily\color{thmredRawSienna!70!black}, bodyfont=\normalfont,
mdframed={
  linewidth=2pt,
  rightline=false, leftline=false, topline=false, bottomline=false,
}
]{thmredname}

\declaretheoremstyle[
  headfont=\bfseries\sffamily\color{RawSienna!70!black}, bodyfont=\normalfont,
  numbered=no,
  mdframed={
    linewidth=2pt,
    rightline=false, topline=false, bottomline=false,
    linecolor=RawSienna, backgroundcolor=RawSienna!1,
  },
  qed=\qedsymbol
]{thmproofbox}

\declaretheoremstyle[
  headfont=\bfseries\sffamily\color{NavyBlue!70!black}, bodyfont=\normalfont,
  numbered=no,
  mdframed={
    linewidth=2pt,
    rightline=false, topline=false, bottomline=false,
    linecolor=NavyBlue, backgroundcolor=NavyBlue!1,
  },
]{thmexplanationbox}

\captionsetup{justification=centering,margin=0.5cm}
\captionsetup{labelfont={color=black}}

\declaretheorem[style=thmgreenbox, name=Definition]{definition}

\declaretheorem[style=thmredbox, name=Proposition]{prop}
\declaretheorem[style=thmredbox, name=Theorem]{theorem}
\declaretheorem[style=thmredbox, name=Lemma]{lemma}
\declaretheorem[style=thmredbox, numbered=no, name=Corollary]{corollary}
\declaretheorem[style=thmredbox, numbered=no, name=Identity]{identity}
\declaretheorem[style=thmredlinesupdown, name=Exercise]{exc}

\declaretheorem[style=thmbluebox, numbered=no, name=Example]{example}
\declaretheorem[style=thmblueline, numbered=no, name=Remark]{remark}

\declaretheorem[style=thmproofbox, name=Proof]{replacementproof}
\renewenvironment{proof}[1][\proofname]{\vspace{-10pt}\begin{replacementproof}}{\end{replacementproof}}

\declaretheorem[style=thmexplanationbox, name=Proof]{tmpexplanation}
\newenvironment{explanation}[1][]{\vspace{-10pt}\begin{tmpexplanation}}{\end{tmpexplanation}}

\newtheorem*{notation}{Notation}
\newtheorem*{note}{Note}
\newtheorem*{previouslyseen}{As previously seen}
\newtheorem*{problem}{Problem}
\newtheorem*{observe}{Observe}
\newtheorem*{property}{Property}
\newtheorem*{intuition}{Intuition}
